\sscp{Possible Austronesian incursion(s): from northeast to southwest}{15-06-03}
This scenario was proposed by \citet{Blust1993}
and clarified in \citet[98]{Blust2008},
“As noted in \citet{Blust1993},
Central Malayo-Polynesian almost certainly was the product of
rapid movement from northeast Indonesia through the central
and southern Moluccas and into the Lesser Sunda islands,
forming a dialect chain that favoured the ready diffusion of
innovations across incipient language boundaries,
and hence destroying any tree-like structure that may have previously
arisen.” 
Since \citet{Blust2012b} rejects a scenario that has movement
from Sulawesi into Maluku (see previous section),
we understand “northeast Indonesia” to mean that Blust sees that
people speaking an Austronesian language or dialect chain moved
from the Halmahera-Bird’s Head region into central Maluku,
through southern Maluku, and travelled westward on into the Lesser Sundas.
No linguistic or archaeological evidence is given to support
such a scenario, so we can only assume it is based on a subgrouping hypothesis
that sees the break-up of PCEMP into CMP
and EMP occurring in or around “northeast Indonesia”.
Yet the break-up of “PCEMP” from the region of north Maluku resulting
in the development of Oceanic languages in connection with the
Lapita complex, and eventually out to Polynesia now appears to have
waning support from archaeology \citep[4]{Bellwood2019b}.
\citet{Bellwood2019b} does not address archaeological
evidence for penetrations southward into Maluku.

\citet[142]{Kamholz2014} also adds important information into
this discussion of north Maluku positing,
“Locating the Proto-SHWNG homeland in southern Cenderawasih Bay
neatly accounts for the spread of the three primary branches
of Proto-SHWNG that are not found in Cenderawasih Bay.
[{\ldots}] Speakers of the third branch (Proto-Raja Ampat-South Halmahera) went westward
along the coast, leaving Cenderawasih Bay and going around the Bird’s Head peninsula.
The homeland of Proto-Raja Ampat-South Halmahera was most likely in the Raja Ampat
archipelago, as three of its five primary branches are found there
(Fiawat, Proto-Ambel-Biga, and Proto-Ma{\Q}ya-Matbat).
Of the remaining two primary branches,
As would have spread eastward to the Bird’s Head,
while Proto-South Halmahera would have spread westward to Halmahera.”

This scenario not only helps explain some things,
but also raises interesting questions.
By having the Austronesian languages of Raja Ampat
and South Halmahera coming from the east,
from southern Cenderawasih Bay,
that helps explain the striking differences
and discontinuities in the historical phonology,
lexicon, syllable structure and grammar between SHWNG languages
and languages in the \tsc{Sula-Buru},
\tsc{Ambon-Seram}, and to a lesser degree \tsc{Seram-Tanimbar-Bomberai} subgroups.
On the surface there are very few linguistic similarities between
these groups and SHWNG languages not accounted for by contact with Papuan languages.
This raises an obvious question: why is there apparently no linguistic
trace evidence of the precursors of “CMP” or an earlier wave
of Austronesian languages in north Maluku (between the Philippines
and central Maluku)? One would think that the Austronesians who
had apparently settled every other island group along their southward
journey, and supposedly had such a far-reaching impact settling the rest
of Maluku and the Lesser Sundas,
would not have completely by-passed such a large group of islands
that had shown themselves to be habitable by humans since around
35,000 BP, and included meat sources such as the tree wallaby
and the cuscus \citep{Bellwood2019a},
as well as fish.
Leapfrogging southward and by-passing north Maluku is certainly
a possibility, but one is still left with the likelihood that different groups
of boats went to different areas,
probably at different times,
and became the precursors of different subgroups.
There was not one monolithic wave-of-advance throughout the region.

Additional concerns about this scenario as the main, or only, incursion
by Austronesian-speaking peoples into Maluku
and the Lesser Sundas include the evidence presented in chapters
\ref{ch:LayAdd}, \ref{ch:OthSub}--\ref{ch:MorLes},
and are reinforced from archaeological dating currently
available for the region (see \citealp{Spriggs2011,Bellwood2019ed},
and discussion in \srf{02-01} and \srf{02-02}).
Archaeological dating has artifacts associated with Austronesian-speaking
peoples present on Timor around 3,800 BP,
much earlier than for north Maluku.
Archaeological dating has artifacts associated with Austronesian-speaking
peoples present in north Maluku near Halmahera (around 3,300
BP on Kayoa), much later than for Timor.
Archaeological dating has artifacts associated with Austronesian-speaking
peoples present in Banda in central Maluku around 3,500 BP.
While 500 years is often trivial in archaeological time,
in this case it coincides with a period of rapid expansion by
Austronesian-speaking peoples into eastern Indonesia
and into the Pacific.
More complete archaeological information in the future may refine
this picture, but that is what is currently available.

So if the incursion from the Halmahera region into Maluku
and the Lesser Sundas is seen as the only incursion,
or the main incursion,
then there is very little linguistic or archaeological evidence
to support it, and significant problems to overcome.
However, if an incursion by Austronesian-speaking peoples via
the Halmahera region is seen as one of several incursions into
the region, then it can be seen as possible,
even though it seems to be based on a hypothesis,
rather than hard data.

For example, there is evidence of retentions such as *p = \it{p},
*b = \it{b}, *ŋ = \it{ŋ},
*k = \it{k}, *-uy = \it{-uy},
*y = \it{y}, along with distinct reflexes of the apicals *d,
*z, *R, *l, *j in \tsc{Sula-Buru}, and retentions of *p = {\#}p,
*b = \it{b}, *ə = **ə,
and *-aw = **-aw in \tsc{Ambon-Seram},
among other things.
Although some of these proto-phonemes in \tsc{Ambon-Seram} are retained
in very few languages
and sometimes only in lexically specific ways.
The diversity and complexity of phonological features,
and structural/typological features attributed to contact in
\tsc{Ambon-Seram} languages suggest a stronger (Papuan) substrate influence there.
But it is possible that people speaking Austronesian languages
not very different from PMP travelling south from the Philippines
by-passed north Maluku and reached central Maluku.
But the different handling of the nasals,
labials, apicals and *ə makes it very difficult to see \tsc{Sula-Buru}
and \tsc{Ambon-Seram} having a shared history through their phonologies.

Additionally, the very different patterns
and discontinuities in the labials,
the apicals, the vowels,
and the vowel-glide sequences,
as well as in the lexicon along with non-cognate morphology in
the subject proclitics, make it very unlikely that \tsc{Timor-Babar},
or \tsc{Bima-Lembata} were derived from or through \tsc{Sula-Buru} or \tsc{Ambon-Seram}.
The opposite scenario, that \tsc{Sula-Buru} or \tsc{Ambon-Seram}
were derived from or through \tsc{Timor-Babar} or \tsc{Bima-Lembata},
is equally unlikely for the same reasons.
The direction of sound changes
and sound correspondences do not match up
and are difficult to reconcile.